% ==========================================================================
% GENERATED by scripts/split_paper.py. Do not edit this file.
%
% Source: imported_papers/grokking/iclr2027_conference_v2.tex
% Part 15: Long run stability and rule switching
%
% Edits belong upstream, in the published paper. To change what the
% chapter includes, change imports/grokking.toml and re-run
% the splitter.
% ==========================================================================
\section{Long run stability and rule switching}
\label{app:long}
If the concentrated circuit were a stop on the way rather than a destination, backpropagation would drift toward the distributed one given a long enough run, and the knob would only be a shortcut. That reading has theory behind it, since the knob's solution is the one max-margin theory calls asymptotically optimal \citep{morwani2024feature} and it sits at backpropagation's own margin normalized norm (Appendix~\ref{app:norm}).

Backpropagation does not drift. Trained for $200{,}000$ epochs, ten times the budget used anywhere else in the paper, its participation settles by roughly $40{,}000$ epochs and is then constant to three decimals, ending at $0.605\,(0.012)$ against the knob's $0.990\,(0.003)$ (Table~\ref{tab:long}). The concentrated circuit is the solution backpropagation occupies, not a waypoint, so selection is the correct description.

The same run shows when credit and circuit come to agree. Both start near uniform, since at initialization every unit receives comparable credit, then fall together through the grokking transition and settle together, staying within about $0.06$ of each other at every checkpoint. Credit participation describes the solution a rule is building as it builds it, which is why we call the agreement a correspondence.

A circuit stable under the rule that built it could still be stable in two ways: in itself, inherited from the grokking transition, or held in place by the rule. A switch separates them, since an inherited circuit would stay put whatever rule runs afterwards and a maintained one should move. We grok at one setting for $20{,}000$ epochs, continue for $30{,}000$ at the other, and run both same rule continuations as drift controls (Table~\ref{tab:long}).

The circuit is maintained. Turning the knob on after backpropagation has grokked moves participation from $0.613\,(0.016)$ to $0.991\,(0.000)$ within the first $5000$ epochs, at test accuracy $0.997$ throughout, while the backpropagation control holds at $0.610\,(0.012)$, so the reorganization is the rule and not the extra training. Removing the knob re-concentrates the distributed circuit to $0.722\,(0.016)$, still falling at the end of the phase, while its control holds at $0.990\,(0.003)$. Both networks had finished the task before the switch and stay at the ceiling through it, so nothing is being learned: the same function is rebuilt over a different set of units because a different rule now assigns the credit. Together with the long run, a rule holds its own circuit indefinitely and moves the other rule's circuit when applied to it.

\begin{table}[!ht]
\caption{The two ends are different solutions rather than one reached at two speeds, and each rule holds its own circuit indefinitely while moving the other rule's circuit when applied to it. Long run stability and rule switching ($p=31$, three seeds, standard deviations in parentheses). In the held block the rule never changes and the two circuit columns are participation at $20{,}000$ and at $200{,}000$ epochs, ten times the budget used anywhere else, with backpropagation's participation constant from about $40{,}000$ epochs onward. In the switched block phase one runs $20{,}000$ epochs and phase two a further $30{,}000$ at the setting named, so the two circuit columns are participation after each phase, and the same rule rows are the drift controls. Credit is the applied per-unit credit participation during phase two, which confirms each phase applied the rule it names, and targeted deletion was measured only on the held runs.}
\label{tab:long}
\begin{center}
\footnotesize
\setlength{\tabcolsep}{4pt}
\begin{tabular}{lcccccc}
\toprule
Condition & Circuit before & Circuit after & Credit & Test & Random del & Targeted del \\
\midrule
\multicolumn{7}{l}{\emph{Rule held, $20{,}000$ to $200{,}000$ epochs}} \\
Backpropagation ($\beta{=}0$) & $0.60$ & $0.605\,(0.012)$ & --- & $0.997$ & $0.673$ & $0.298$ \\
Equalized ($\beta{=}1$) & $0.99$ & $0.990\,(0.003)$ & --- & $0.975$ & $0.775$ & $0.599$ \\
\midrule
\multicolumn{7}{l}{\emph{Rule switched, phase one to phase two}} \\
$\beta{=}0 \to \beta{=}1$ & $0.613\,(0.016)$ & $0.991\,(0.000)$ & $0.999$ & $0.997$ & $0.800$ & --- \\
$\beta{=}1 \to \beta{=}0$ & $0.990\,(0.003)$ & $0.722\,(0.016)$ & $0.719$ & $0.990$ & $0.703$ & --- \\
$\beta{=}0 \to \beta{=}0$ & $0.613\,(0.016)$ & $0.610\,(0.012)$ & $0.601$ & $0.997$ & $0.675$ & --- \\
$\beta{=}1 \to \beta{=}1$ & $0.990\,(0.003)$ & $0.990\,(0.003)$ & $1.000$ & $0.975$ & $0.775$ & --- \\
\bottomrule
\end{tabular}
\end{center}
\end{table}
